% 引言
\section{Introduction}

% 低比特权重量化能够降低大语言模型的存储需求与权重读取开销。向量量化（VQ）通过一个索引联合表示多个参数，相比逐标量编码具有更灵活的表示能力。近年来，QTIP 和 LLVQ 等方法进一步利用 trellis 或格结构提高有效量化维度，但也带来相应的搜索与部署约束：无结构高维码本的近邻搜索和缓存开销随维度快速增长；trellis、网格和格点方法则通过路径或格约束降低搜索难度，同时减少子块独立查表与局部改写的灵活性。因此，本文关注一个核心问题：在保留低维码字可查表性和局部可改写性的前提下，扩大跨位置的联合选择范围，能否使模型质量接近甚至达到高维结构化量化方法的水平？
Low-bit weight quantization reduces the storage and weight-loading costs of large language models. Vector quantization (VQ) represents multiple parameters jointly with one index and therefore offers greater representational flexibility than scalar quantization~\cite{grayneuhoff,gershogray}. Recent methods such as QTIP and LLVQ increase the effective quantization dimension using trellis or lattice structures~\cite{qtip,llvq}, but they also introduce search and deployment constraints: nearest-neighbor search and cache costs grow rapidly for unstructured high-dimensional codebooks, while trellises and lattices reduce search complexity through path or lattice constraints at the cost of independent lookup and local rewritability. This paper therefore asks whether expanding joint selection across positions can approach the quality of high-dimensional structured quantization while retaining lookup-friendly, locally rewritable low-dimensional codewords.

% 回答这一问题，需要区分向量内部的联合表示与多个向量之间的联合选择。以共享尺度 VQ 为例，同一尺度组内的多个位置分别选择低维码字，但共用同一个尺度参数。由于最优尺度取决于整组码字的组合，它又会反过来改变各位置候选的相对优劣。于是，即便每次单位置更新后都重新优化尺度，逐点方法仍可能陷入联合停滞：没有任何单位置更新能够继续降低误差，但同时改变多个位置仍可能获得真实下降。图 1 用两阶段轨迹展示这一现象；右图进一步给出被接受更新中索引改动数的分布。本文将这一现象形式化为共享尺度下的联合选择问题，研究在不提高基础码字维度的情况下，如何以可控成本发现并验证这类联合更新。
Answering this question requires distinguishing joint representation within a vector from joint selection across multiple vectors. In shared-scale VQ, positions in the same scale group select low-dimensional codewords independently but share one scale. Because the optimal scale depends on the complete codeword combination and in turn changes the relative ordering of candidates, a pointwise method may stagnate even when it re-optimizes the scale after every single-position update: no individual change reduces the error, whereas a coordinated multi-position change can. Figure~\ref{fig:motivation} illustrates this behavior with two-stage trajectories and reports the number of indices changed by accepted updates. We formulate this phenomenon as joint selection under a shared scale and study how to discover and verify such updates at controlled cost without increasing the base-codeword dimension.

\begin{figure}[t]
\centering
\includegraphics[width=\linewidth]{figure1_motivation_code.png}
\caption{
% （a）三个共享尺度组在单位置索引更新与尺度重优化阶段逐步进入平台；切换到联合索引更新后，归一化真实层输出误差继续下降。每条曲线均以切换时的误差归一化为 1。（b）被接受更新的索引改动数分布；多索引更新占多数。
(a) Three shared-scale groups approach a plateau under single-index updates with scale re-optimization; after switching to joint-index updates, the normalized true layer-output error continues to decrease. Each curve is normalized to 1 at the switch. (b) Distribution of the number of changed indices among accepted updates; multi-index updates are the majority.
}
\label{fig:motivation}
\end{figure}

% 为实现可控的联合更新，我们提出共享尺度下的码字协同优化方法。候选池由若干探针尺度下的局部短名单取并集，使搜索不绑定在当前尺度的候选排序上。随后利用完整校准残差建立相切可分离上界，使每个位置的候选代价成为共享尺度的二次函数。利用正尺度域内公共项可消去的性质，我们将候选比较转化为直线下包络，并合并各位置的切换点。在每个代表区间中，联合码字组合保持不变，只需求解一个受合法尺度约束的一维二次问题。该过程直接更新硬索引与实际尺度，不对索引引入概率参数或反向传播；其最优性限定于当前候选池中的代理目标。
To obtain controllable joint updates, we propose cooperative codeword optimization under a shared scale. Candidate pools are formed by taking the union of local shortlists at several probe scales, so the search is not tied to the candidate ordering at the current scale. We then use the full calibration residual to construct a tangent separable upper bound, under which each candidate cost is quadratic in the shared scale. By canceling the common term over the positive-scale domain, candidate comparison becomes a lower-envelope problem over lines. After merging switch points across positions, the joint codeword combination is fixed within each representative interval, leaving only a one-dimensional quadratic problem over legal scales. The procedure directly updates hard indices and stored scales without probabilistic index parameters or backpropagation; its optimality is restricted to the surrogate objective over the current candidate pools.

% 实验结果验证了联合决策的有效性。在 Qwen3-4B 和 Qwen3-8B 上，尺度与码字分配联合优化将 WikiText-2 PPL 分别从 22.69/14.41 降至 12.18/9.76，明显优于仅优化共享尺度。在 Qwen3.8-27B 上，本方法相比 QTIP 和 LLVQ 取得更低的 WikiText-2/C4 PPL，并在 LiveCodeBench v6 和 GPQA-D 上获得更高分数。维度消融进一步表明，在约 2 bpw 的实际码率下，低维基础码字结合跨位置联合优化能够取得有竞争力的模型质量。
Experiments support the effectiveness of joint decisions. On Qwen3-4B and Qwen3-8B, jointly optimizing the scale and codeword assignments reduces WikiText-2 perplexity from 22.69/14.41 to 12.18/9.76, substantially outperforming scale-only optimization. On Qwen3.8-27B, our method achieves lower WikiText-2/C4 perplexity than QTIP and LLVQ and higher scores on LiveCodeBench v6 and GPQA-D. The dimension ablation further shows that low-dimensional base codewords combined with cross-position joint optimization provide competitive model quality at an actual rate of approximately 2 bpw. Figure~\ref{fig:qwen3-4b-overview} visualizes the Qwen3-4B results.

\begin{figure}[t]
\centering
\includegraphics[width=\linewidth]{figure2_qwen3_4b_ppl_mean.png}
\caption{
% Qwen3-4B 结果概览。（a）WikiText-2 与 C4 PPL 的等权平均值，越低越好。（b）所选方法在五个下游任务上的结果；完整数值比较以表格为准。
Overview of Qwen3-4B results. (a) Equal-weight mean of WikiText-2 and C4 perplexity; lower is better. (b) Results of selected methods on five downstream tasks. The tables provide the complete numerical comparison.
}
\label{fig:qwen3-4b-overview}
\end{figure}

% 本文将基础码字维度与联合优化范围分开设计，贡献归纳为三点：
Our contributions are threefold:
\begin{itemize}
\item
% 提出固定低维表示下的共享尺度联合优化方法，用多尺度探针并集生成候选，并通过完整残差感知的评价协调多个硬码字与尺度。
We propose shared-scale joint optimization under a fixed low-dimensional representation, using a union of multi-scale probe candidates and full-residual-aware evaluation to coordinate multiple hard codewords and their shared scale.
\item
% 建立相切可分离上界及下包络区间求解，给出候选池内的上界最优性以及相应的理论边界。
We construct a tangent separable upper bound and a lower-envelope interval solver, and establish upper-bound optimality within the candidate pools together with its theoretical limits.
\item
% 通过同起点协同消融、基础维度消融和多模型量化评测，考察联合选择对模型质量与优化效率的影响。
We evaluate the effect of joint selection on model quality and optimization cost through same-start ablations, a base-dimension ablation, and multi-model quantization experiments.
\end{itemize}
